% Options for packages loaded elsewhere
\PassOptionsToPackage{unicode}{hyperref}
\PassOptionsToPackage{hyphens}{url}
\documentclass[
]{article}
\usepackage{xcolor}
\usepackage[margin=1in]{geometry}
\usepackage{amsmath,amssymb}
\setcounter{secnumdepth}{-\maxdimen} % remove section numbering
\usepackage{iftex}
\ifPDFTeX
  \usepackage[T1]{fontenc}
  \usepackage[utf8]{inputenc}
  \usepackage{textcomp} % provide euro and other symbols
\else % if luatex or xetex
  \usepackage{unicode-math} % this also loads fontspec
  \defaultfontfeatures{Scale=MatchLowercase}
  \defaultfontfeatures[\rmfamily]{Ligatures=TeX,Scale=1}
\fi
\usepackage{lmodern}
\ifPDFTeX\else
  % xetex/luatex font selection
\fi
% Use upquote if available, for straight quotes in verbatim environments
\IfFileExists{upquote.sty}{\usepackage{upquote}}{}
\IfFileExists{microtype.sty}{% use microtype if available
  \usepackage[]{microtype}
  \UseMicrotypeSet[protrusion]{basicmath} % disable protrusion for tt fonts
}{}
\makeatletter
\@ifundefined{KOMAClassName}{% if non-KOMA class
  \IfFileExists{parskip.sty}{%
    \usepackage{parskip}
  }{% else
    \setlength{\parindent}{0pt}
    \setlength{\parskip}{6pt plus 2pt minus 1pt}}
}{% if KOMA class
  \KOMAoptions{parskip=half}}
\makeatother
\usepackage{longtable,booktabs,array}
\usepackage{calc} % for calculating minipage widths
% Correct order of tables after \paragraph or \subparagraph
\usepackage{etoolbox}
\makeatletter
\patchcmd\longtable{\par}{\if@noskipsec\mbox{}\fi\par}{}{}
\makeatother
% Allow footnotes in longtable head/foot
\IfFileExists{footnotehyper.sty}{\usepackage{footnotehyper}}{\usepackage{footnote}}
\makesavenoteenv{longtable}
\usepackage{graphicx}
\makeatletter
\newsavebox\pandoc@box
\newcommand*\pandocbounded[1]{% scales image to fit in text height/width
  \sbox\pandoc@box{#1}%
  \Gscale@div\@tempa{\textheight}{\dimexpr\ht\pandoc@box+\dp\pandoc@box\relax}%
  \Gscale@div\@tempb{\linewidth}{\wd\pandoc@box}%
  \ifdim\@tempb\p@<\@tempa\p@\let\@tempa\@tempb\fi% select the smaller of both
  \ifdim\@tempa\p@<\p@\scalebox{\@tempa}{\usebox\pandoc@box}%
  \else\usebox{\pandoc@box}%
  \fi%
}
% Set default figure placement to htbp
\def\fps@figure{htbp}
\makeatother
\setlength{\emergencystretch}{3em} % prevent overfull lines
\providecommand{\tightlist}{%
  \setlength{\itemsep}{0pt}\setlength{\parskip}{0pt}}
\usepackage{bookmark}
\IfFileExists{xurl.sty}{\usepackage{xurl}}{} % add URL line breaks if available
\urlstyle{same}
\hypersetup{
  hidelinks,
  pdfcreator={LaTeX via pandoc}}

\author{}
\date{\vspace{-2.5em}}

\begin{document}

\section{EIG and ELLR -- Theory \& Battleship
Simulator}\label{eig-and-ellr-theory-battleship-simulator}

\begin{center}\rule{0.5\linewidth}{0.5pt}\end{center}

\subsection{Prompt 1}\label{prompt-1}

\begin{quote}
See \href{https://arxiv.org/pdf/2510.20886}{``SHOOT FIRST, ASK QUESTIONS
LATER? BUILDING RATIONAL AGENTS THAT EXPLORE AND ACT LIKE PEOPLE''}.
They use expected information gain, EIG(qt \textbar{} x, H1:t)=EIGε(qt
\textbar{} x, H1:t) = Hb(ε + (1 − 2ε) pt) − Hb(ε), to choose question qt
at time t given the partially observed board x and history H1:t of all
resulting answers of questions and results of actions. Where Hb is the
binary entropy. Show mathematically how EIG differs from using the
expected log-likelihood ratio log(LR(H=s;E=a˜t \textbar qt, x, H1:t)),
where hypothesis H=board s and evidence E=answer to qt is a˜t. And the
expectation is taken over all boards s∈S⊢x where S⊢x = \{ s ∈ S : s is
consistent with x \}, i.e.~all boards consistent with partially observed
board x, with probabiity πt(s) = Pr(S = s \textbar{} x, H1:t).
LR(H=s;E=a˜t \textbar qt, x, H1:t) = Pr(E=a˜t\textbar qt, s, x, H1:t) /
Pr(E=a˜t\textbar qt, not(s), x, H1:t). So log-likelihood ratio is log
Pr(E=a˜t\textbar qt, s, x, H1:t) - log Pr(E=a˜t\textbar qt, not(s), x,
H1:t). And not(s) is all boards other than s.
\end{quote}

\subsection{Response, Claude Sonnet
4.6}\label{response-claude-sonnet-4.6}

This is a clean information-theoretic question I can answer directly.
Here's the full derivation.

\begin{center}\rule{0.5\linewidth}{0.5pt}\end{center}

\subsection{Setup and Notation}\label{setup-and-notation}

Fix question \(q_t\), board \(x\), and history \(H_{1:t}\). Write
\(\pi_s \equiv \pi_t(s)\) for brevity, and let
\(p_s(\tilde{a}) \equiv \Pr(\tilde{A}_t = \tilde{a} \mid s, q_t)\).
Under the BSC(\(\varepsilon\)) noise model,

\[p_s(1) = \varepsilon + (1-2\varepsilon)\,\mathbf{1}\{f_q(s)=1\}, \qquad p_s(0) = 1 - p_s(1).\]

The \textbf{marginal} answer distribution (over all boards) is

\[\bar{p}(\tilde{a}) \equiv \Pr(\tilde{A}_t = \tilde{a} \mid q_t, x, H_{1:t}) = \sum_{s'} \pi_{s'}\, p_{s'}(\tilde{a}),\]

and the \textbf{leave-\(s\)-out} marginal (over all boards except \(s\))
is

\[\bar{p}_{-s}(\tilde{a}) \equiv \Pr(\tilde{A}_t = \tilde{a} \mid q_t, \text{not}(s), x, H_{1:t}) = \frac{\bar{p}(\tilde{a}) - \pi_s\, p_s(\tilde{a})}{1 - \pi_s}.\]

\begin{center}\rule{0.5\linewidth}{0.5pt}\end{center}

\subsection{Both Quantities as Weighted KL
Divergences}\label{both-quantities-as-weighted-kl-divergences}

\textbf{EIG} is the mutual information
\(I(S;\tilde{A}_t \mid x, H_{1:t})\), which expands as

\[\boxed{\mathrm{EIG} = \sum_s \pi_s\; D_{\mathrm{KL}}\!\bigl(p_s \;\|\; \bar{p}\bigr)}\]

\textbf{Expected log-LR} is, by definition,

\[\mathrm{E}[\log \mathrm{LR}] = \sum_s \pi_s \sum_{\tilde{a}} p_s(\tilde{a})\;\log\frac{p_s(\tilde{a})}{\bar{p}_{-s}(\tilde{a})} = \sum_s \pi_s\; D_{\mathrm{KL}}\!\bigl(p_s \;\|\; \bar{p}_{-s}\bigr).\]

\[\boxed{\mathrm{E}[\log \mathrm{LR}] = \sum_s \pi_s\; D_{\mathrm{KL}}\!\bigl(p_s \;\|\; \bar{p}_{-s}\bigr)}\]

The \textbf{only structural difference} is the reference distribution
inside the KL divergence: EIG references \(\bar{p}\), the mixture over
\emph{all} boards; E{[}log LR{]} references \(\bar{p}_{-s}\), the
mixture \emph{leaving out} board \(s\) itself.

\begin{center}\rule{0.5\linewidth}{0.5pt}\end{center}

\subsection{Exact Relationship}\label{exact-relationship}

Because
\(\bar{p}(\tilde{a}) = \pi_s\, p_s(\tilde{a}) + (1-\pi_s)\,\bar{p}_{-s}(\tilde{a})\),
the ratio of the two reference distributions is

\[\frac{\bar{p}(\tilde{a})}{\bar{p}_{-s}(\tilde{a})} = (1-\pi_s) + \pi_s\,\underbrace{\frac{p_s(\tilde{a})}{\bar{p}_{-s}(\tilde{a})}}_{\displaystyle\equiv\, \mathrm{LR}_s(\tilde{a})}.\]

Therefore

\[D_{\mathrm{KL}}(p_s \| \bar{p}_{-s}) - D_{\mathrm{KL}}(p_s \| \bar{p}) = \sum_{\tilde{a}} p_s(\tilde{a})\,\log\frac{\bar{p}(\tilde{a})}{\bar{p}_{-s}(\tilde{a})} = \mathbb{E}_{\tilde{a}\sim p_s}\!\left[\log\bigl(1-\pi_s + \pi_s\,\mathrm{LR}_s(\tilde{A})\bigr)\right].\]

Summing over \(s\) with weights \(\pi_s\):

\[\boxed{\mathrm{E}[\log \mathrm{LR}] - \mathrm{EIG} = \sum_s \pi_s\;\mathbb{E}_{\tilde{a}\sim p_s}\!\left[\log\bigl(1 - \pi_s + \pi_s\,\mathrm{LR}_s(\tilde{A})\bigr)\right].}\]

\begin{center}\rule{0.5\linewidth}{0.5pt}\end{center}

\subsection{\texorpdfstring{EIG \(\leq\) E{[}log LR{]}
Always}{EIG \textbackslash leq E{[}log LR{]} Always}}\label{eig-leq-elog-lr-always}

Since \(\mathrm{KL}(P\|\cdot)\) is \textbf{convex} in its second
argument, and \(\bar{p}\) is the convex mixture

\[\bar{p} = \pi_s\, p_s + (1-\pi_s)\,\bar{p}_{-s},\]

we get

\[D_{\mathrm{KL}}(p_s \| \bar{p}) \leq \pi_s\, D_{\mathrm{KL}}(p_s \| p_s) + (1-\pi_s)\, D_{\mathrm{KL}}(p_s \| \bar{p}_{-s}) = (1-\pi_s)\, D_{\mathrm{KL}}(p_s \| \bar{p}_{-s}).\]

Hence, term-by-term and after taking the weighted sum:

\[\mathrm{EIG} \;\leq\; \sum_s \pi_s(1-\pi_s)\, D_{\mathrm{KL}}(p_s \| \bar{p}_{-s}) \;\leq\; \mathrm{E}[\log \mathrm{LR}].\]

EIG is always a \textbf{lower bound} on the expected log-LR, with the
gap controlled by the self-weight \(\pi_s\).

\begin{center}\rule{0.5\linewidth}{0.5pt}\end{center}

\subsection{When They Coincide}\label{when-they-coincide}

Taylor-expand the correction term for small \(\pi_s\):

\[\log(1 - \pi_s + \pi_s\,\mathrm{LR}_s) = \pi_s(\mathrm{LR}_s - 1) - \tfrac{\pi_s^2}{2}(\mathrm{LR}_s-1)^2 + O(\pi_s^3).\]

Taking the expectation over \(\tilde{a} \sim p_s\) and noting
\(\mathbb{E}_{p_s}[\mathrm{LR}_s] = \sum_{\tilde a} p_s(\tilde a)\,p_s(\tilde a)/\bar p_{-s}(\tilde a) \geq 1\):

\[\mathrm{E}[\log \mathrm{LR}] - \mathrm{EIG} = O\!\left(\sum_s \pi_s^2\right).\]

For a uniform prior over \(N\) boards, \(\pi_s = 1/N\), so the gap is
\(O(1/N)\). As \(N\to\infty\) (or more generally as the belief \(\pi_t\)
becomes diffuse), \textbf{the two measures coincide}. With the SMC
particle approximation used in the paper, this occurs as the number of
particles grows.

\begin{center}\rule{0.5\linewidth}{0.5pt}\end{center}

\subsection{Intuitive Summary}\label{intuitive-summary}

\begin{longtable}[]{@{}
  >{\raggedright\arraybackslash}p{(\linewidth - 4\tabcolsep) * \real{0.3333}}
  >{\raggedright\arraybackslash}p{(\linewidth - 4\tabcolsep) * \real{0.3333}}
  >{\raggedright\arraybackslash}p{(\linewidth - 4\tabcolsep) * \real{0.3333}}@{}}
\toprule\noalign{}
\begin{minipage}[b]{\linewidth}\raggedright
\end{minipage} & \begin{minipage}[b]{\linewidth}\raggedright
Reference distribution
\end{minipage} & \begin{minipage}[b]{\linewidth}\raggedright
Board \(s\) included?
\end{minipage} \\
\midrule\noalign{}
\endhead
\bottomrule\noalign{}
\endlastfoot
\textbf{EIG} & \(\bar p = \sum_{s'}\pi_{s'} p_{s'}\) & Yes --- \(s\)
inflates its own answer probability \\
\textbf{E{[}log LR{]}} &
\(\bar p_{-s} = (\bar p - \pi_s p_s)/(1-\pi_s)\) & No ---
leave-one-out \\
\end{longtable}

The log-LR measures how ``surprising'' the observed answer is
\emph{relative to every other hypothesis}, making it strictly larger.
EIG instead compares each board to the \textbf{population-average}
answer distribution, which is already partially ``contaminated'' by that
board's own contribution. This self-contamination shrinks the apparent
divergence, giving EIG \(\leq\) E{[}log LR{]}, with equality only when
no single board dominates the belief (i.e., \(\max_s \pi_s \to 0\)).

\begin{center}\rule{0.5\linewidth}{0.5pt}\end{center}

\subsection{Prompt 2}\label{prompt-2}

\begin{quote}
given the difference, can an argument be made to prefer either one over
the other?
\end{quote}

\subsection{Response, Claude Sonnet
4.6}\label{response-claude-sonnet-4.6-1}

This is actually a conceptual question I can address directly from first
principles --- no web research needed here.

There are principled arguments on both sides, and the right choice
depends on what the question-asker cares about.

\begin{center}\rule{0.5\linewidth}{0.5pt}\end{center}

\subsection{Arguments for EIG}\label{arguments-for-eig}

\textbf{1. Coherent Bayesian accounting.} EIG = \(I(S;\tilde{A}_t)\) is
the \emph{expected reduction in Shannon entropy of the full belief}
\(\pi_t\). This means it accounts for the agent's \textbf{global
epistemic state}: a question is good if it compresses the entire
hypothesis space, not just if it discriminates one board from the rest.
This is exactly what Bayesian Experimental Design (BED) optimizes for,
and what the paper formalizes.

\textbf{2. Additivity over turns.} Because entropy is additive, EIG has
a clean sequential interpretation: total information gained across turns
is the sum of per-turn EIGs (modulo correlations from belief updates).
This compositionality is very convenient for planning and for comparing
strategies across games of different lengths.

\textbf{3. Answer-symmetric.} EIG is symmetric in
\(\tilde{a} \in \{0, 1\}\): a question is maximally informative when
\(p_t \approx 1/2\), i.e., when ``yes'' and ``no'' are equally likely.
The log-LR is not symmetric in this way --- it depends on which answer
is observed, and the expectation over answers introduces asymmetries
depending on the prior.

\textbf{4. Practical computability.} As shown in the paper (Eq. 4), EIG
under BSC(\(\varepsilon\)) has a \textbf{closed form} in terms of
\(p_t\), which is itself just a particle-weight average. This makes it
tractable in the SMC approximation. The expected log-LR requires a
leave-one-out marginal \(\bar{p}_{-s}\) for every particle, which is
\(O(N)\) more expensive to compute.

\textbf{5. Robustness to degenerate beliefs.} When \(\pi_s \to 1\) for
some board \(s\) (near-certainty), the log-LR diverges because
\(\bar{p}_{-s}\) can be very different from \(p_s\), making the ratio
extreme. EIG stays bounded in \([0, 1-H_b(\varepsilon)]\) regardless.

\begin{center}\rule{0.5\linewidth}{0.5pt}\end{center}

\subsection{Arguments for E{[}log LR{]}}\label{arguments-for-elog-lr}

\textbf{1. Hypothesis-discriminating interpretation.} The log-LR
directly measures the \textbf{weight of evidence} (in the
Bayesian/Peircean sense) that the answer provides \emph{for a specific
board \(s\) against all alternatives}. If the goal is to confirm or rule
out particular hypotheses --- as in medical diagnosis or targeted
scientific experiments --- this is more semantically natural.

\textbf{2. Connection to hypothesis testing.} The log-LR is the central
quantity in Neyman-Pearson testing and sequential probability ratio
tests (SPRT). If the agent's downstream task is a \textbf{binary
decision} (e.g., ``is the ship at this location?'') rather than full
belief compression, optimizing expected log-LR aligns with minimizing
expected type-I/II error more directly than EIG does.

\textbf{3. Penalizes overconfident priors less harshly.} Because
\(\bar{p}_{-s}\) excludes \(s\)'s self-contribution, it measures
discriminability against the \emph{rest of the world} rather than
against a mixture that already includes \(s\). This means the log-LR
does not artificially deflate the value of a question just because the
agent already assigns high probability to the correct board --- it
rewards questions that sharply separate the true board from alternatives
even when the prior is concentrated.

\textbf{4. Closer to causal/interventional reasoning.} The log-LR can be
interpreted as: ``if I run this experiment and observe answer
\(\tilde{a}\), by how much should I update my log-odds on board \(s\)?''
This maps cleanly onto practical scientific reasoning --- researchers
often think in terms of evidence strength for specific theories, not
global entropy reduction.

\begin{center}\rule{0.5\linewidth}{0.5pt}\end{center}

\subsection{The Deeper Tension}\label{the-deeper-tension}

The two criteria reflect a \textbf{fundamental trade-off} in active
learning and experimental design:

\begin{longtable}[]{@{}
  >{\raggedright\arraybackslash}p{(\linewidth - 4\tabcolsep) * \real{0.3333}}
  >{\raggedright\arraybackslash}p{(\linewidth - 4\tabcolsep) * \real{0.3333}}
  >{\raggedright\arraybackslash}p{(\linewidth - 4\tabcolsep) * \real{0.3333}}@{}}
\toprule\noalign{}
\begin{minipage}[b]{\linewidth}\raggedright
\end{minipage} & \begin{minipage}[b]{\linewidth}\raggedright
EIG
\end{minipage} & \begin{minipage}[b]{\linewidth}\raggedright
E{[}log LR{]}
\end{minipage} \\
\midrule\noalign{}
\endhead
\bottomrule\noalign{}
\endlastfoot
\textbf{Goal} & Compress the whole belief & Discriminate specific
hypotheses \\
\textbf{Framework} & Information theory (Shannon) & Evidence theory
(Bayes factors) \\
\textbf{Scales like} & \(\sum_s \pi_s \log(\pi_s / \bar p)\) &
\(\sum_s \pi_s \log(\pi_s / \bar p_{-s})\) \\
\textbf{Coincide when} & Belief is diffuse (\(\pi_s \ll 1\)) & --- \\
\textbf{Preferred when} & Many hypotheses, uncertain prior & Few
hypotheses, or decision-focused \\
\end{longtable}

The gap between them,
\(\sum_s \pi_s \mathbb{E}[\log(1 - \pi_s + \pi_s \mathrm{LR}_s)]\), is
largest when the belief is \textbf{concentrated} --- late in the game
after many questions. This is precisely when the two criteria give the
most different advice about what question to ask next.

In the Battleship setting specifically, EIG is the right choice for the
reasons the paper emphasizes: the hypothesis space is huge, the prior
starts diffuse, and the goal is efficient global belief compression. But
in a setting like diagnostic medicine --- where you have a short list of
candidate diseases and a patient must be treated --- the expected log-LR
(or equivalently, optimizing Bayes factors) is arguably the more
principled criterion, because the clinician cares about discrimination,
not entropy.

\begin{center}\rule{0.5\linewidth}{0.5pt}\end{center}

\subsection{Prompt 3}\label{prompt-3}

\begin{quote}
is there a scenario that you can represent in code and simulate that
clearly illustrates the different path in queries given the same initial
setting? Say, a simple board search problem: much smaller board than in
paper? If so, before coding it, generate a detailed explanation of the
scenario.
\end{quote}

\subsection{Response, Claude Sonnet
4.6}\label{response-claude-sonnet-4.6-2}

This is a conceptual design question I can work through from our
mathematical analysis --- no research needed. Let me lay out the
scenario carefully before any code is written.

\begin{center}\rule{0.5\linewidth}{0.5pt}\end{center}

\subsection{The Scenario}\label{the-scenario}

\subsubsection{Grid and Hypothesis
Space}\label{grid-and-hypothesis-space}

A \textbf{4×4 grid} with a single \textbf{length-2 ship} (domino). The
full hypothesis space has 24 positions (12 horizontal, 12 vertical).
After some initial observations --- say, confirmed misses at several
cells --- the belief has been narrowed to exactly \textbf{4 remaining
hypotheses}:

\begin{longtable}[]{@{}llll@{}}
\toprule\noalign{}
ID & Ship position & Cells occupied & Weight \(\pi_s\) \\
\midrule\noalign{}
\endhead
\bottomrule\noalign{}
\endlastfoot
s1 & Horizontal, row 0 & (0,0) and (0,1) & \textbf{0.45} \\
s2 & Horizontal, row 1 & (1,0) and (1,1) & \textbf{0.45} \\
s3 & Vertical, col 2 & (0,2) and (1,2) & 0.05 \\
s4 & Vertical, col 2 & (1,2) and (2,2) & 0.05 \\
\end{longtable}

The asymmetric prior reflects that earlier hits or revealed water tiles
made the two horizontal positions in the top-left highly probable,
leaving a small residual probability for two vertical positions in
column 2. This is the critical state: belief is heavily concentrated on
two equally plausible dominant boards, with two low-probability
alternatives still in play.

\subsubsection{The Two Candidate
Questions}\label{the-two-candidate-questions}

From our analysis, the divergence between EIG and E{[}log LR{]} is
sharpest when the belief is \textbf{concentrated and asymmetric with
respect to the question structure}. The most revealing pair of candidate
questions here is:

\textbf{Q\_cell}: \emph{``Is cell (0,1) occupied?''}

\begin{itemize}
\tightlist
\item
  s1 → Yes (cell (0,1) is the right half of the horizontal ship in row
  0)
\item
  s2, s3, s4 → No
\item
  Marginal answer probability:
  \(p_t \approx 0.45 \times 0.9 + 0.55 \times 0.1 = 0.46\) (slightly
  below 0.5)
\end{itemize}

\textbf{Q\_row}: \emph{``Is there any ship in row 0?''}

\begin{itemize}
\tightlist
\item
  s1 → Yes (occupies cells in row 0)
\item
  s3 → Yes (has a cell at (0,2))
\item
  s2, s4 → No
\item
  Marginal answer probability:
  \(p_t \approx (0.45 + 0.05) \times 0.9 + 0.50 \times 0.1 = 0.50\)
  (exactly balanced)
\end{itemize}

\subsubsection{Why EIG and E{[}log LR{]} Disagree
Here}\label{why-eig-and-elog-lr-disagree-here}

\textbf{EIG} measures
\(H_b(\varepsilon + (1-2\varepsilon)p_t) - H_b(\varepsilon)\), which is
maximized at \(p_t = 0.5\). Q\_row achieves exactly this by combining s1
(weight 0.45) and s3 (weight 0.05) into the ``Yes'' partition, perfectly
balancing the expected answer distribution. Q\_cell falls just short,
with \(p_t = 0.46\). So \textbf{EIG prefers Q\_row}.

\textbf{E{[}log LR{]}} computes
\(\sum_s \pi_s \log \frac{p_s(\tilde{a})}{\bar{p}_{-s}(\tilde{a})}\).
For Q\_cell and board s1:
\[\bar{p}_{-s1}(\text{Yes}) = \frac{\bar{p}(\text{Yes}) - \pi_{s1} \cdot p_{s1}(\text{Yes})}{1 - \pi_{s1}} = \frac{0.46 - 0.45 \times 0.9}{0.55} = \frac{0.055}{0.55} = 0.10\]

So the log-LR for s1 under Q\_cell is
\(\log(0.9/0.10) = \log 9 \approx 2.20\) --- nearly perfect
discrimination. For Q\_row, s3 ``dilutes'' the Yes partition:
\(\bar{p}_{-s1}(\text{Yes})\) is now larger (the leave-s1-out marginal
is no longer near zero), so the log-LR for s1 drops to about 1.65. The
45\% weight on s1 makes this difference decisive: \textbf{E{[}log LR{]}
prefers Q\_cell}.

The divergence in numbers (with \(\varepsilon = 0.1\)):

\begin{longtable}[]{@{}llll@{}}
\toprule\noalign{}
Question & \(p_t\) & EIG (bits) & E{[}log LR{]} \\
\midrule\noalign{}
\endhead
\bottomrule\noalign{}
\endlastfoot
Q\_row & 0.50 & \textbf{0.531} & 1.548 \\
Q\_cell & 0.46 & 0.529 & \textbf{1.628} \\
\end{longtable}

The gap in EIG is tiny (0.002 bits). The gap in E{[}log LR{]} is
meaningful (0.080). They rank oppositely.

\subsubsection{The Structural Reason for
Disagreement}\label{the-structural-reason-for-disagreement}

Q\_row achieves its perfect \(p_t = 0.5\) by \textbf{bundling s1 with
s3} in the Yes partition. This is informationally efficient in a global
sense --- the answer compresses the hypothesis space maximally. But if
the answer is Yes, you still cannot distinguish s1 (the dominant board,
\(\pi = 0.45\)) from s3 (\(\pi = 0.05\)). The posterior on s1 after a
Yes answer is \(\approx 0.81\), not 1.0.

Q\_cell \textbf{perfectly isolates s1}. A Yes answer concentrates
posterior mass on s1 at \(\approx 0.88\). A No answer elevates s2 to
\(\approx 0.75\) while redistributing small weight to s3/s4
proportionally. In neither case does a minority board ``contaminate''
the signal about the dominant hypothesis.

E{[}log LR{]} rewards this because it explicitly weights the
log-likelihood ratio for each board by that board's prior probability.
The 2.20 log-LR for s1 (weight 0.45) dominates the sum, making Q\_cell
win despite its slightly worse global entropy reduction.

\subsubsection{What the Simulation Should
Show}\label{what-the-simulation-should-show}

The simulation runs many full identification games from this belief
state under each strategy, greedy over one criterion at each step:

\begin{itemize}
\tightlist
\item
  \textbf{Performance when true board is s1 or s2} (90\% of games):
  Q\_cell-based strategy identifies the board faster, because the first
  question provides extreme discrimination of the dominant board
\item
  \textbf{Performance when true board is s3 or s4} (10\% of games):
  Q\_row-based strategy does slightly better here, because the row query
  gives useful information about the column-2 vertical ships too
\item
  \textbf{Overall expected questions to identification}: E{[}log LR{]}
  strategy wins overall despite the rare-case disadvantage, because the
  90\% dominant case dominates the expectation
\end{itemize}

The simulation should also show the \textbf{belief trajectory} --- how
posterior entropy evolves step by step under each strategy --- which
will visually illustrate that EIG compresses entropy uniformly while
E{[}log LR{]} collapses the dominant hypothesis faster at the cost of
slower convergence in the tail.

\begin{center}\rule{0.5\linewidth}{0.5pt}\end{center}

Should I proceed to code this up now?

\begin{center}\rule{0.5\linewidth}{0.5pt}\end{center}

\subsection{Prompt 4}\label{prompt-4}

\begin{quote}
yes. make it clear what candidate questions are available to pick from.
i suspect Q\_cell for each cell, Q\_row for each row and Q\_column for
each column, right?
\end{quote}

\subsection{Response, Claude Sonnet
4.6}\label{response-claude-sonnet-4.6-3}

Yes, exactly right --- cell questions (one per cell = 16), row questions
(4), column questions (4), giving 24 total. Let me code it up.Here's
what the simulator contains and what to look for.

\begin{center}\rule{0.5\linewidth}{0.5pt}\end{center}

\subsubsection{Question catalogue (24
total)}\label{question-catalogue-24-total}

\begin{longtable}[]{@{}lll@{}}
\toprule\noalign{}
Type & Count & Example \\
\midrule\noalign{}
\endhead
\bottomrule\noalign{}
\endlastfoot
\textbf{Cell} & 16 & \texttt{Cell\ (0,1)} --- ``Is (row 0, col 1)
occupied?'' \\
\textbf{Row} & 4 & \texttt{Row\ 0} --- ``Is any ship in row 0?'' \\
\textbf{Col} & 4 & \texttt{Col\ 0} --- ``Is any ship in column 0?'' \\
\end{longtable}

Many are uninformative at the initial prior (e.g., \texttt{Col\ 3} ---
no hypothesis has a ship there). The checkbox ``Include uninformative''
reveals them; by default only the 11 informative questions are shown.

\begin{center}\rule{0.5\linewidth}{0.5pt}\end{center}

\subsubsection{What the ranking table shows at the initial belief
state}\label{what-the-ranking-table-shows-at-the-initial-belief-state}

\textbf{EIG \#1 → \texttt{Row\ 0}} (P(Yes) = 0.50 exactly; perfectly
balanced partition \{H1, H3\} vs \{H2, H4\})

\textbf{E{[}log LR{]} \#1 → \texttt{Cell\ (0,0)}} (and Cell(0,1),
Cell(1,0), Cell(1,1) all tie) --- these perfectly isolate the dominant
hypothesis (H1 or H2, weight 0.45) from all others, yielding a very high
leave-one-out KL divergence despite P(Yes) = 0.46 rather than 0.50.

The \texttt{Rank\ Δ} column shows how dramatically the two criteria
re-rank each question --- \texttt{Row\ 0} will be EIG \#1 but ELLR \#3
or \#4, while the cell questions are EIG \#2 but ELLR \#1.

\begin{center}\rule{0.5\linewidth}{0.5pt}\end{center}

\subsubsection{How to use it}\label{how-to-use-it}

\begin{enumerate}
\def\labelenumi{\arabic{enumi}.}
\tightlist
\item
  \textbf{Click a hypothesis} to designate it as ground truth for
  interactive play
\item
  \textbf{``Ask EIG-best''} or \textbf{``Ask E{[}log LR{]}-best''} ---
  the belief updates after a noisy answer, and the question rankings
  re-sort live
\item
  \textbf{Watch the divergence banner} --- it turns yellow the moment
  the criteria disagree, green when they agree (tends to happen once the
  belief is already concentrated)
\item
  \textbf{``Simulate 500×4 games''} --- runs the full greedy strategy
  for each criterion on each true hypothesis and reports average
  questions-to-convergence and accuracy
\end{enumerate}

\begin{center}\rule{0.5\linewidth}{0.5pt}\end{center}

\ldots{}

\ldots{}

\ldots{}

\begin{center}\rule{0.5\linewidth}{0.5pt}\end{center}

\subsection{Prompt 14}\label{prompt-14}

\begin{quote}
based upon theory, is it wrong-headed to think of the ELLR metric as
adopting a 0-1 loss function, thus seeking max a posteriori choices
vs.~the EIG metric as adopting a least-squares loss function, thus
seeking a maximum expectation over the posterior? Seems like there's at
least an analogous thinking between these. As such, it makes sense to me
that for an explicitly enumerated set of hypotheses and tractable
computation, for Battleship -- in which you're rewarded for hits and not
for near misses -- that ELLR is preferred. But if the hypotheses are too
many to enumerate all of them -- i.e., high-dimensional; large
Battleship board -- that EIG might be preferred. Another way to say is
that when there is lots of uncertainty, either due to lack of info or
approximate computation or high epsilon noise parameter, then we opt for
exploring as guided by EIG; but when the uncertainty becomes better
resolved and posterior mass concentrates on few enough hypotheses to be
explicitly computable, then we'd switch to exploiting as guided by ELLR.
Do any of these ideas hold any water? Or have I just taken things too
far based on near-similarities that don't hold up under careful
theoretical computation?
\end{quote}

\subsection{Response, Claude Sonnet
4.6}\label{response-claude-sonnet-4.6-4}

This question is conceptual and draws directly on our prior derivations
--- I can reason through it directly without research.

Your intuitions are genuinely interesting and some of them hold up well,
but others don't quite survive careful theoretical inspection. Let me go
through each thread.

\begin{center}\rule{0.5\linewidth}{0.5pt}\end{center}

\subsection{Thread 1: ELLR as 0-1 loss / MAP, EIG as squared-error loss
/ posterior
mean}\label{thread-1-ellr-as-0-1-loss-map-eig-as-squared-error-loss-posterior-mean}

This analogy is \emph{suggestive but not precise}. Here's what's
actually happening:

\textbf{EIG} optimizes for the question that maximally compresses the
posterior --- it's agnostic about what you'll \emph{do} with the
information. The downstream decision rule is not specified by EIG
itself. You could follow EIG question-selection with a MAP decision, a
posterior-mean decision, or anything else. EIG is purely an epistemic
criterion.

\textbf{ELLR} is also purely an epistemic criterion, not a decision
criterion. It measures how well the observed answer discriminates the
true hypothesis from all alternatives, weighted by prior. It doesn't
directly encode a loss function either.

So neither criterion \emph{is} a loss function. However, your analogy
points at something real: ELLR is more sensitive to \textbf{what the
dominant hypothesis is} (because it weights discrimination by
\(\pi_s\)), while EIG is sensitive to \textbf{the shape of the whole
distribution}. In that sense, ELLR behaves more like a criterion that
``cares about identifying the winner,'' which is spiritually closer to
MAP reasoning. EIG behaves more like a criterion that ``cares about the
whole distribution,'' which is spiritually closer to posterior-mean
reasoning.

The precise connection is: ELLR is the expected \textbf{log posterior
odds} update --- after observing \(\tilde{a}\), the log-odds on board
\(s\) shift by \(\log \frac{p_s(\tilde{a})}{\bar{p}_{-s}(\tilde{a})}\),
and ELLR is the expectation of this shift over answers and boards. So
ELLR is directly the expected gain in \textbf{discriminative evidence}
in the Bayes-factor sense, which is the right criterion if your goal is
to concentrate posterior mass on one hypothesis as fast as possible.
That is indeed structurally MAP-adjacent.

\begin{center}\rule{0.5\linewidth}{0.5pt}\end{center}

\subsection{Thread 2: Battleship rewards hits not near-misses, so ELLR
is
preferred}\label{thread-2-battleship-rewards-hits-not-near-misses-so-ellr-is-preferred}

This is the sharpest and most defensible claim you've made, and I think
it's essentially correct. The argument runs:

In Battleship the reward structure is \textbf{binary per cell}: you
either hit or miss, and a miss at distance 1 from a ship is worth
exactly the same as a miss at distance 10. This is precisely a 0-1 loss
on the move decision. Given that, the right move criterion is MAP (fire
at the highest-probability cell), and --- as we saw in the paper with
MBayes --- explicitly marginalizing over the posterior to pick the MAP
cell is exactly what works.

The question-selection criterion should then be chosen to \emph{serve}
that downstream MAP move decision. ELLR, by rewarding questions that
sharply separate the dominant hypothesis from all others, is more
directly aligned with driving the posterior toward a state where the MAP
cell is reliably correct. EIG, by seeking balanced partitions, can
``waste'' a question finely splitting probability mass between two
hypotheses that both point to the same region of the board --- the
entropy goes down but the MAP cell doesn't change.

This is a genuine and non-trivial insight. It parallels results in the
Bayesian experimental design literature on the distinction between
\textbf{entropy-reduction} designs (minimize posterior entropy) and
\textbf{decision-theoretic} designs (minimize expected loss under the
posterior given a specific action space). When your action space is
discrete and your loss is 0-1, the decision-theoretic design and the
entropy-reduction design can diverge, and the decision-theoretic one is
preferable.

\begin{center}\rule{0.5\linewidth}{0.5pt}\end{center}

\subsection{Thread 3: EIG preferred in high-dimensional / intractable
settings; ELLR preferred when posterior is concentrated and
tractable}\label{thread-3-eig-preferred-in-high-dimensional-intractable-settings-ellr-preferred-when-posterior-is-concentrated-and-tractable}

This is where I'd push back most carefully. The argument has two
separable parts:

\textbf{Tractability:} It's true that ELLR requires computing a
leave-one-out marginal \(\bar{p}_{-s}\) for every particle, making it
\(O(N)\) more expensive than EIG. So in high-dimensional settings where
you're running SMC with many particles, EIG is cheaper. This is a
practical point, not a theoretical one.

\textbf{High uncertainty → EIG, concentrated posterior → ELLR:} This is
more subtle and I don't think it holds up cleanly. The gap between EIG
and ELLR is
\(\sum_s \pi_s \mathbb{E}[\log(1 - \pi_s + \pi_s \mathrm{LR}_s)]\),
which is largest when the posterior is \textbf{concentrated} (large
\(\pi_s\) for some \(s\)). So ELLR and EIG diverge \emph{most} when the
posterior is already concentrated --- meaning ELLR is most distinctively
useful precisely when you already have a good idea of the answer.

When the posterior is diffuse (early in the game, high uncertainty), EIG
\(\approx\) ELLR to \(O(1/N)\), so the choice barely matters. This
somewhat inverts your framing: it's not that you switch \emph{to} ELLR
as uncertainty resolves --- rather, ELLR is always better in principle
for the 0-1 reward setting, but the difference only becomes practically
significant as the posterior concentrates.

\textbf{The explore/exploit framing:} This is appealing but I'd restate
it slightly. EIG is not really an exploration criterion and ELLR is not
really an exploitation criterion --- both are question-selection
criteria that always ask questions (exploring). The exploit step is the
\emph{move decision} (firing at a cell), which is separate. What you're
actually describing is the transition from asking questions to taking
shots, which in the paper is handled by DBayes (the one-step lookahead
decision). That's orthogonal to the EIG vs ELLR choice for question
selection.

\begin{center}\rule{0.5\linewidth}{0.5pt}\end{center}

\subsection{Summary}\label{summary}

\begin{longtable}[]{@{}
  >{\raggedright\arraybackslash}p{(\linewidth - 2\tabcolsep) * \real{0.5000}}
  >{\raggedright\arraybackslash}p{(\linewidth - 2\tabcolsep) * \real{0.5000}}@{}}
\toprule\noalign{}
\begin{minipage}[b]{\linewidth}\raggedright
Your claim
\end{minipage} & \begin{minipage}[b]{\linewidth}\raggedright
Verdict
\end{minipage} \\
\midrule\noalign{}
\endhead
\bottomrule\noalign{}
\endlastfoot
ELLR is MAP-adjacent, EIG is posterior-mean-adjacent & Approximately
right, but neither is strictly a loss function \\
Battleship's 0-1 reward makes ELLR the better question criterion & Solid
--- this is a genuine decision-theoretic argument \\
EIG preferred when hypothesis space is too large to enumerate &
Partially right as a \emph{tractability} argument; not a theoretical
superiority \\
High uncertainty → EIG, concentrated posterior → ELLR & Inverted: the
gap between them is largest (ELLR most distinctively useful) when
posterior is \emph{already} concentrated \\
Explore with EIG, exploit with ELLR & The exploit step is the move
decision, not question selection --- so this conflates two separate
choices \\
\end{longtable}

The deepest correct insight in your thinking is the second one: given a
0-1 reward structure, ELLR is the more decision-theoretically coherent
question criterion. The rest are useful analogies that need some
refinement to hold precisely.

\begin{center}\rule{0.5\linewidth}{0.5pt}\end{center}

\subsection{Prompt 15}\label{prompt-15}

\begin{quote}
good explanations -- though in thread 1, i didn't mean to imply that the
metrics themselves were loss functions; i meant what you've called out
in thread 2, that search guided by ellr aligns more closely when you
have a 0-1 loss function than does search with eig. good clarification
on thread 3. thanks, you've helped to refine my thinking on this topic
of experimental searches \& active learning which i've studied for a
long time.
\end{quote}

\subsection{Response, Claude Sonnet
4.6}\label{response-claude-sonnet-4.6-5}

Thank you --- and I understood what you meant in thread 1 on reflection.
The rephrasing is cleaner: EIG aligns with searching to minimize a
squared-error-like loss on the posterior (reduce uncertainty uniformly),
while ELLR aligns with searching to minimize 0-1 loss on the final
identification decision. That's a crisp and defensible distinction.

What strikes me most about your framing is that it connects experimental
design to the classic statistical question of what your estimator is
\emph{for}. The choice of search criterion is really a choice of
implicit downstream loss, and the mismatch between criterion and loss is
what creates the divergence we've been simulating --- EIG wastes
question budget achieving balanced partitions that don't move the MAP
cell, exactly because it isn't ``aware'' that the game scores hits not
entropy reductions.

The explore/exploit thread is also worth keeping in a refined form: not
EIG vs ELLR as explore vs exploit, but rather the \emph{question-asking
phase itself} as exploration, and the transition from asking to shooting
(DBayes) as the exploit decision. The paper's result that DBayes +
ELLR-style move selection (MBayes) outperforms pure EIG-question
selection alone is consistent with this --- the two mechanisms operate
at different levels of the decision hierarchy.

If you ever revisit this in a research context, the natural next
question seems to be whether there's a unified criterion that
interpolates between EIG and ELLR as a function of posterior
concentration --- essentially adapting the question-selection criterion
dynamically as the game progresses, switching from entropy-reduction to
discriminative-evidence mode as mass concentrates. That gap appears to
be open in the literature.

\begin{center}\rule{0.5\linewidth}{0.5pt}\end{center}

\subsection{Prompt 16}\label{prompt-16}

\begin{quote}
oh, and what if epsilon swaps p, i.e.~low signal-to-noise ratio in
answers or conversily epsilon goes to zero: what are impacts on metric
choice?
\end{quote}

\subsection{Response, Claude Sonnet
4.6}\label{response-claude-sonnet-4.6-6}

This is a clean theoretical question that falls directly out of our
earlier derivations --- I can work through it directly.

\begin{center}\rule{0.5\linewidth}{0.5pt}\end{center}

\subsection{What ε does to both
metrics}\label{what-ux3b5-does-to-both-metrics}

Recall the BSC(ε) model: the noisy answer \(\tilde{A}\) flips the true
answer with probability ε. The signal-to-noise ratio of a single answer
is essentially \(1 - 2\varepsilon\), which runs from 1 (perfect) at ε=0
to 0 (pure noise) at ε=½.

Both metrics depend on ε through the same channel parameter, but they
respond differently.

\begin{center}\rule{0.5\linewidth}{0.5pt}\end{center}

\subsection{EIG as a function of ε}\label{eig-as-a-function-of-ux3b5}

\[\mathrm{EIG}_\varepsilon(q) = H_b\!\left(\varepsilon + (1-2\varepsilon)\,p_t\right) - H_b(\varepsilon)\]

The two terms behave as follows. \(H_b(\varepsilon)\) is the entropy
injected by the channel itself --- it rises from 0 at ε=0 to 1 at ε=½.
The first term \(H_b(\varepsilon + (1-2\varepsilon)p_t)\) is the total
entropy of the noisy answer. Their difference is the mutual information
that survives the channel.

\textbf{At ε → 0:} EIG → \(H_b(p_t)\), the binary entropy of the
noiseless marginal. The metric works exactly as designed --- a perfectly
balanced question with \(p_t = ½\) achieves the maximum of 1 bit.

\textbf{At ε → ½:} Both terms approach \(H_b(½) = 1\) bit, so EIG → 0
for all questions regardless of \(p_t\). The channel is pure noise; no
question is worth asking. EIG correctly reflects that no information can
be extracted.

\textbf{Critically:} the \emph{ranking} of questions by EIG is preserved
under ε. Since EIG is a monotone function of \(|p_t - ½|\) for any fixed
ε (questions closer to \(p_t = ½\) always rank higher), and ε doesn't
change which question achieves the most balanced partition, \textbf{ε
scales the absolute EIG values uniformly downward but does not change
the rank ordering of questions}. EIG's choice of best question is
ε-invariant.

\begin{center}\rule{0.5\linewidth}{0.5pt}\end{center}

\subsection{E{[}log LR{]} as a function of
ε}\label{elog-lr-as-a-function-of-ux3b5}

\[\mathrm{E}[\log\mathrm{LR}] = \sum_s \pi_s\, D_{\mathrm{KL}}(p_s \| \bar{p}_{-s})\]

where \(p_s(\tilde{a}) \in \{\varepsilon,\, 1-\varepsilon\}\) depending
on whether hypothesis \(s\) answers Yes or No to \(q\).

\textbf{At ε → 0:} \(p_s(\tilde{a}) \in \{0, 1\}\) --- answers are
deterministic. For a hypothesis \(s\) that answers Yes,
\(p_s(\text{Yes}) = 1\) and \(p_s(\text{No}) = 0\), so the KL divergence
from \(\bar{p}_{-s}\) becomes \(\log(1/\bar{p}_{-s}(\text{Yes}))\),
which can be very large if the leave-one-out marginal assigns low
probability to Yes. The metric becomes extremely sensitive to questions
that would produce answers inconsistent with all other hypotheses --- it
strongly rewards questions that \emph{perfectly isolate} individual
hypotheses. The discrimination-maximizing character of ELLR is sharpest
here.

\textbf{At ε → ½:} \(p_s(\tilde{a}) → ½\) for all hypotheses and all
questions --- the noisy answer distribution becomes identical for all
\(s\), collapsing to Bernoulli(½). Therefore \$p\_s
\to \bar\{p\}\_\{-s\} \to \$ Bernoulli(½) for all \(s\), and every KL
divergence goes to zero. So E{[}log LR{]} → 0 uniformly, just like EIG.
But the path to zero is different.

\textbf{The key asymmetry:} Unlike EIG, the \emph{ranking} of questions
by ELLR \textbf{is sensitive to ε}. As ε increases, the contrast between
\(p_s(\text{Yes}) = 1-\varepsilon\) and \(p_s(\text{No}) = \varepsilon\)
shrinks. This compresses the differences in KL divergence across
questions. Questions that were highly discriminating at low ε (large KL
for the dominant hypothesis) shrink faster than questions with moderate
discrimination, because the KL between two distributions close to each
other in \(\{ε, 1-ε\}\) depends on the ratio
\(\frac{1-\varepsilon}{\varepsilon}\), which is very sensitive to ε.
Specifically:

\[D_{\mathrm{KL}}(p_s \| \bar{p}_{-s}) \propto (1-2\varepsilon)^2 \cdot \text{(noiseless discrimination)}\]

to first order in small \((1-2\varepsilon)\), meaning ELLR shrinks as
the \emph{square} of signal strength whereas EIG shrinks roughly
linearly in mutual information. So \textbf{ELLR degrades faster than EIG
as ε increases}.

\begin{center}\rule{0.5\linewidth}{0.5pt}\end{center}

\subsection{The crossover implication for metric
choice}\label{the-crossover-implication-for-metric-choice}

Putting it together:

\begin{longtable}[]{@{}
  >{\raggedright\arraybackslash}p{(\linewidth - 6\tabcolsep) * \real{0.2500}}
  >{\raggedright\arraybackslash}p{(\linewidth - 6\tabcolsep) * \real{0.2500}}
  >{\raggedright\arraybackslash}p{(\linewidth - 6\tabcolsep) * \real{0.2500}}
  >{\raggedright\arraybackslash}p{(\linewidth - 6\tabcolsep) * \real{0.2500}}@{}}
\toprule\noalign{}
\begin{minipage}[b]{\linewidth}\raggedright
Regime
\end{minipage} & \begin{minipage}[b]{\linewidth}\raggedright
EIG ranking
\end{minipage} & \begin{minipage}[b]{\linewidth}\raggedright
ELLR ranking
\end{minipage} & \begin{minipage}[b]{\linewidth}\raggedright
Preferred
\end{minipage} \\
\midrule\noalign{}
\endhead
\bottomrule\noalign{}
\endlastfoot
ε → 0 (perfect SNR) & stable, picks \(p_t = ½\) question & most
discriminating, picks dominant-isolating question & \textbf{ELLR} ---
discrimination maximally sharp \\
ε moderate & stable & compressed but still distinct from EIG &
\textbf{ELLR} still, for 0-1 reward, but gap narrows \\
ε → ½ (pure noise) & all → 0, uniform & all → 0, but faster &
\textbf{EIG} --- degrades more gracefully; ELLR collapses first \\
\end{longtable}

The crossover is real and has a clean interpretation. At low ε, answers
are trustworthy, so it makes sense to ask the question that most sharply
discriminates your best hypothesis --- you'll act on the answer. At high
ε, answers are unreliable, so the question ranking should reflect
\emph{robustness to noise}, which is better captured by EIG's
balanced-partition criterion. A balanced question with \(p_t = ½\) is
maximally robust to noise because flipping the answer still gives you
useful information (just in the other direction), whereas a highly
discriminating question with \(p_t \approx 1-\varepsilon\) for the
dominant hypothesis becomes nearly useless when noise can flip it.

\begin{center}\rule{0.5\linewidth}{0.5pt}\end{center}

\subsection{Connection to your earlier
insight}\label{connection-to-your-earlier-insight}

This reinforces your thread 2 point in a new way. The 0-1 loss argument
for ELLR assumed answers are reliable enough to act on. As ε grows, that
assumption weakens: the answer is no longer a clean signal about which
hypothesis to bet on, so the MAP-adjacent reasoning behind ELLR loses
its force, and the entropy-reduction reasoning behind EIG --- which
never assumed you'd act immediately on any single answer --- becomes
relatively more valuable. In the limit of pure noise, the only rational
strategy is to ask balanced questions and aggregate many weak signals,
exactly what EIG prescribes.

So ε can be read as a \emph{trust parameter}: trust the answers and use
ELLR; distrust them and use EIG.

\begin{center}\rule{0.5\linewidth}{0.5pt}\end{center}

\end{document}
